\documentclass[aspectratio=169, 10pt]{beamer}

% ==============================================================================
% 1. PACKAGES & FONTS
% ==============================================================================
\usepackage[utf8]{inputenc}
\usepackage[T1]{fontenc}
\usepackage{helvet} % Matches the poster's font choice 
\usepackage[english]{babel}
\usepackage{booktabs} % For nice tables
\usepackage{tikz}
\usepackage{tcolorbox}
\usepackage{graphicx}
\usepackage{multicol}

% Bibliography setup
\usepackage[style=numeric, backend=biber]{biblatex}
\addbibresource{../../soc/literatura.bib}

% Set default font family to sans-serif (Helvetica)
\renewcommand{\familydefault}{\sfdefault}

% ==============================================================================
% 2. COLOR PALETTE (Extracted from poster.tex )
% ==============================================================================
\definecolor{NTCPurple}{RGB}{151, 13, 157}   % Primary accent
\definecolor{NTCPurple2}{RGB}{179, 94, 184}  % Secondary accent
\definecolor{NTCLight}{RGB}{231, 210, 233}   % Light tint
\definecolor{NTCGrey}{RGB}{60, 60, 60}       % Body text
\definecolor{LightGrey}{RGB}{240, 240, 240}  % Backgrounds

% ==============================================================================
% 3. BEAMER THEME CUSTOMIZATION
% ==============================================================================

% Remove navigation symbols
\setbeamertemplate{navigation symbols}{}

% Color mappings
\setbeamercolor{structure}{fg=NTCPurple}
\setbeamercolor{normal text}{fg=NTCGrey}
\setbeamercolor{title}{fg=NTCPurple}
\setbeamercolor{subtitle}{fg=NTCPurple2}
\setbeamercolor{frametitle}{fg=NTCPurple}
\setbeamercolor{section in toc}{fg=NTCPurple}

% Itemize bullets (Square/Circle to match scientific style)
\setbeamertemplate{itemize item}{\scriptsize\raise1.25pt\hbox{\color{NTCPurple}$\blacksquare$}}
\setbeamertemplate{itemize subitem}{\scriptsize\raise1.25pt\hbox{\color{NTCPurple2}$\blacktriangleright$}}
\setbeamertemplate{itemize subsubitem}{\tiny\raise1.5pt\hbox{\color{NTCGrey}$\bullet$}}

% Custom Block Style (Matches the poster's tcolorbox )
\setbeamercolor{block title}{bg=NTCPurple2, fg=white}
\setbeamercolor{block body}{bg=NTCLight!35, fg=NTCGrey}
\setbeamertemplate{blocks}[rounded][shadow=false]

% Footer (Simple with page number)
\setbeamertemplate{footline}{
  \hbox{
  \begin{beamercolorbox}[wd=\paperwidth,ht=2.25ex,dp=1ex,right]{date in head/foot}
    \usebeamerfont{date in head/foot}\insertshortauthor\hspace*{2em}
    \insertframenumber{} / \inserttotalframenumber\hspace*{2ex} 
  \end{beamercolorbox}}
  \vskip0pt
}

% Header (Horizontal line in NTCPurple2)
\setbeamertemplate{headline}{
  \vskip10pt
  \hspace{0.5cm}\textcolor{NTCPurple2}{\rule{0.95\paperwidth}{2pt}}
}

% ==============================================================================
% 4. METADATA
% ==============================================================================
\title{Hybrid Quantum-Classical Neural Network\\for Pneumonia Detection}
\subtitle{CSASC 2026 Graduation Project}
\author[M. Forgó]{Michal Forgó\inst{1,2} \\ \small{Supervisor: Bc. Jan Boháč}}
\institute{
  \inst{1} Secondary School of Informatics and Financial Services (SŠINFIS)\\
  \inst{2} New Technologies Research Centre (NTC), UWB
}
\date{\today}

% ==============================================================================
% 5. PRESENTATION CONTENT
% ==============================================================================
\begin{document}

% --- TITLE SLIDE ---
{
\setbeamertemplate{headline}{} % Hide header on title slide
\begin{frame}
  \begin{columns}
    \column{0.7\textwidth}
      \vspace{1cm}
      {\Huge \textbf{\textcolor{NTCPurple}{Hybrid Quantum-Classical Neural Network}}\\[0.2cm] 
      {\Large \textcolor{NTCGrey}{Pneumonia Detection in Chest X-Rays}}\\[1cm]
      
      \textbf{Michal Forgó}\\
      \footnotesize \textit{University of West Bohemia \& SŠINFIS}\\[0.2cm]
      
      \textcolor{NTCPurple2}{\rule{5cm}{1pt}}
      
    \column{0.3\textwidth}
      \centering
      % PLACEHOLDERS FOR LOGOS
      \includegraphics[width=0.8\linewidth, height=2cm, keepaspectratio]{example-image-a}\\ \vspace{0.5cm}
      \includegraphics[width=0.8\linewidth, height=2cm, keepaspectratio]{example-image-b}
  \end{columns}
\end{frame}
}

% --- TOC SLIDE ---
\begin{frame}{Overview}
  \tableofcontents
\end{frame}

% --- SECTION 1 ---
\section{Introduction \& Scientific Context}

% Standard text slide
\begin{frame}{Clinical Motivation & Context}
  \begin{itemize}
    \item \textbf{The Problem:} Pneumonia causes $\sim$2.5 million deaths annually[cite: 8].
    \item \textbf{Current State:} Deep Learning (ResNet, DenseNet) achieves high accuracy ($>85\%$) but requires massive parameter counts ($164.8k+$ per layer)[cite: 9].
    \item \textbf{The Opportunity:} NISQ quantum devices offer parameter-efficient learning via VQA[cite: 10].
  \end{itemize}
  
  \vspace{0.5cm}
  
  % Custom block matching the poster style
  \begin{block}{Research Hypothesis}
    A quantum decision layer (VQC with $\sim$100 parameters) can approach classical discrimination while reducing memory footprint, enabling Edge/IoT deployment[cite: 12].
  \end{block}
\end{frame}

% --- SECTION 2 ---
\section{Methodology}

% Two Column Slide with Table
\begin{frame}{Hybrid Architecture Design}
  \textbf{Design Rationale:} Separate feature learning (classical) from decision making (quantum).

  \vspace{0.5cm}

  \begin{columns}[T] % T aligns columns at the top
    \column{0.5\textwidth}
       \textbf{Pipeline Flow:}
       \begin{enumerate}
           \item \textbf{Input:} Chest X-Ray ($224 \times 224$).
           \item \textbf{Feature Extraction:} ConvNeXt-Tiny (frozen) $\rightarrow$ 768-D vector.
           \item \textbf{Compression:} PCA/Autoencoder ($768 \rightarrow 64$).
           \item \textbf{Classification:} VQC (6 qubits, 54 params).
       \end{enumerate}

    \column{0.5\textwidth}
      \centering
      \footnotesize
       \begin{tabular}{llc}
         \toprule
         \textbf{Stage} & \textbf{Operation} & \textbf{Params} \\
         \midrule
         Features & ConvNeXt-Tiny & 0 (frozen) \\
         Reduction & PCA/Autoencoder & 0 (fixed) \\
         \textbf{Decision} & \textbf{VQC (6-qubit)} & \textbf{54} \\
         \bottomrule
       \end{tabular}
      \vspace{0.2cm}
      \captionof{table}{Parameter counts [cite: 22]}
  \end{columns}
\end{frame}

% --- SECTION 3 ---
\section{Results}

% Slide with Math/Equations
\begin{frame}{Quantum Circuit Theory}
  \textbf{Amplitude Encoding:}
  Mapping vector $x \in \mathbb{R}^{64}$ to quantum state:
  $$ |\psi(x)\rangle = \sum_{i=0}^{63} x_i |i\rangle \in \mathcal{H}_{2^6} $$
  
  \vspace{0.3cm}
  
  \textbf{VQC Prediction:}
  Defined by the expectation value:
  $$ f(x,\theta) = \langle \psi(x) | U^\dagger(\theta)\, Z_0\, U(\theta) | \psi(x) \rangle \in [-1,1] $$
  
  \vspace{0.3cm}
  \small{Optimized via weighted mean squared error on simulator (PennyLane)[cite: 43].}
\end{frame}

% Comparison Slide
\begin{frame}{Performance vs. Efficiency}
  The VQC decision head achieves competitive results with massive parameter reduction.

  \vspace{0.5cm}
  
  \begin{columns}
      \column{0.5\textwidth}
      \textbf{Parameter Efficiency:}
      \begin{itemize}
          \item Classical MLP: 2,113 params.
          \item Quantum VQC: \textbf{54 params}.
          \item \textcolor{NTCPurple}{\textbf{Reduction: 39$\times$}}[cite: 58].
      \end{itemize}
      
      \column{0.5\textwidth}
      \textbf{Diagnostic Accuracy (AUC):}
      \begin{itemize}
          \item Classical MLP: 0.940 AUC.
          \item Hybrid VQC: \textbf{0.974 AUC}[cite: 66].
      \end{itemize}
  \end{columns}
  
  \vspace{0.5cm}
  \centering
  \textit{Note: Specificity drops due to dataset shift and barren plateaus[cite: 69].}
\end{frame}

% --- END SLIDE ---
{
\setbeamertemplate{headline}{}
\begin{frame}
  \centering
  \vspace{1cm}
  \huge \textcolor{NTCPurple}{\textbf{Thank You for Your Attention}}
  
  \vspace{1cm}
  \normalsize
  \textbf{Questions?}
  
  \vspace{1.5cm}
  \footnotesize
  \textcolor{NTCGrey}{Contact: mforgo@example.com}
\end{frame}
}

% --- REFERENCES ---
\begin{frame}{References}
  \printbibliography[heading=none]
\end{frame}

\end{document}